\documentclass[../../../include/open-logic-section]{subfiles}

\begin{document}

\iftag{FOL}
      {\olfileid{fol}{seq}{der}}
      {\olfileid{pl}{seq}{der}}

\olsection{推导}

% Part: first-order-logic
% Chapter: sequent-calculus
% Section: derivations

\begin{explain}
我们介绍了初始相继式的样子，也给出了推理规则。相继式演算中的推导就是由这些归纳生成的：每个推导要么本身就是一个初始相继式，要么由一个或两个推导再接一个推理规则构成。
\end{explain}

\begin{defn}[$\Log{LK}$ 推导]
相继式~$S$ 的\emph{$\Log{LK}$-推导}是一棵满足以下条件的有限相继式树：
\begin{enumerate}
\item 树的顶端相继式是初始相继式。
\item 树的底端相继式是~$S$。
\item 树中除~$S$ 之外的每个相继式，都是某个推理规则正确应用的前提，而该规则的结论恰好位于树中该相继式的正下方。
\end{enumerate}
此时，我们称~$S$ 是该推导的\emph{末位相继式}（end-sequent），并称~$S$ 是\emph{可在 $\Log{LK}$ 中推导的}（或称 $\Log{LK}$-可推导的）。
\end{defn}

\begin{ex}
每一个初始相继式，例如 $!C \Sequent !C$，本身就是一个推导。我们可以通过应用某个推理规则，比如 $\LeftR{\Weakening}$ 规则，从此推导得到一个新推导：
\begin{prooftree}
\Axiom$ \Gamma \fCenter \Delta $
\RightLabel{\LeftR{\Weakening}}
\UnaryInf$ !A, \Gamma \fCenter \Delta$
\end{prooftree}
然而，这个规则是通用的：我们可以把规则中的 $!A$ 替换成任何语句，比如换成~$!D$。如果前提匹配我们的初始相继式 $!C \Sequent !C$，那就意味着 $\Gamma$ 和 $\Delta$ 都只是~$!C$，那么结论就会是 $!D, !C \Sequent !C$。因此，下面就是一个推导：
\begin{prooftree}
\Axiom$ !C \fCenter !C $
\RightLabel{\LeftR{\Weakening}}
\UnaryInf$ !D, !C \fCenter !C$
\end{prooftree}
我们现在可以应用另一个规则，比如 $\LeftR{\Exchange}$ 规则，它允许我们交换左边的两个语句。所以，下面也是一个正确的推导：
\begin{prooftree}
\Axiom$ !C \fCenter !C $
\RightLabel{\LeftR{\Weakening}}
\UnaryInf$ !D, !C \fCenter !C$
\RightLabel{\LeftR{\Exchange}}
\UnaryInf$ !C, !D \fCenter !C$
\end{prooftree}
在这次规则应用中，规则原本的形式是
\begin{prooftree}
\Axiom$ \Gamma, !A, !B, \Pi \fCenter \Delta $
\RightLabel{\LeftR{\Exchange}}
\UnaryInf$ \Gamma, !B, !A, \Pi \fCenter \Delta,$
\end{prooftree}
其中 $\Gamma$ 和 $\Pi$ 都为空，$\Delta$ 就是 $!C$，而 $!A$ 和 $!B$ 的角色分别由 $!D$ 和~$!C$ 扮演。以几乎同样的方式，我们也看到
\begin{prooftree}
\Axiom$ !D \fCenter !D $
\RightLabel{\LeftR{\Weakening}}
\UnaryInf$ !C, !D \fCenter !D$
\end{prooftree}
是一个推导。现在我们可以取这两个推导，使用 $\RightR{\land}$ 规则将它们组合起来。该规则是
\begin{prooftree}
\Axiom$\Gamma \fCenter \Delta, !A$
\Axiom$ \Gamma \fCenter \Delta, !B$
\RightLabel{\RightR{\land}}
\BinaryInf$ \Gamma \fCenter \Delta, !A \land !B $
\end{prooftree}
在我们的例子中，前提必须与那两个推导各自的末位相继式相匹配。这意味着 $\Gamma$ 是 $!C, !D$，$\Delta$ 是空集，$!A$ 是 $!C$ 而 $!B$ 是 $!D$。因此，要使该推理正确，结论就应该是 $!C, !D \Sequent !C
\land !D$。
\begin{prooftree}
\Axiom$ !C \fCenter !C $
\RightLabel{\LeftR{\Weakening}}
\UnaryInf$ !D, !C \fCenter !C$
\RightLabel{\LeftR{\Exchange}}
\UnaryInf$ !C, !D \fCenter !C$
\Axiom$ !D \fCenter !D $
\RightLabel{\LeftR{\Weakening}}
\UnaryInf$ !C, !D \fCenter !D$
\RightLabel{\RightR{\land}}
\BinaryInf$ !C, !D \fCenter !C \land !D $
\end{prooftree}
当然，我们也可以调换两个前提的顺序，那样的话 $!A$ 就是 $!D$ 而 $!B$ 就是~$!C$。
\begin{prooftree}
\Axiom$ !D \fCenter !D $
\RightLabel{\LeftR{\Weakening}}
\UnaryInf$ !C, !D \fCenter !D$
\Axiom$ !C \fCenter !C $
\RightLabel{\LeftR{\Weakening}}
\UnaryInf$ !D, !C \fCenter !C$
\RightLabel{\LeftR{\Exchange}}
\UnaryInf$ !C, !D \fCenter !C$
\RightLabel{\RightR{\land}}
\BinaryInf$ !C, !D \fCenter !D \land !C $
\end{prooftree}
\end{ex}

\end{document}